\chapter{{\LaTeX}による本ドキュメントの編集方法}
\section{環境構築}
Dev Containerによって{\LaTeX}の開発環境をまるごとコンテナ化することで、面倒な{\LaTeX}の環境構築を回避しつつ、本ドキュメントを編集することができます。ただし、依然としてコンテナを動かすための環境構築は必須であり、ここではその方法を説明します。手順さえ踏めば、{\LaTeX}のインストールとは違って、謎のエラーに苦しむこともないはずです。

Dev Containerについてのより詳しい解説はVisual Studio Codeのウェブサイト（\url{https://code.visualstudio.com/docs/devcontainers/containers}）を参照してください。ここでの説明はこのサイトの和訳と抜粋にすぎません。

\subsection{Dockerのインストール}
\subsubsection{Windowsにインストールする場合}
Docker Desktopをインストールしてください。WindowsのエディションがHomeの場合はWSL2機能を有効にする必要があります。

\subsubsection{Windows上のWSL2（Ubuntu）にインストールする場合}
\begin{enumerate}
    \item 以下2つをインストールしてください。具体的な方法はそれぞれのウェブサイトを参照してください。
    \begin{itemize}
        \item Docker Engine: \url{https://docs.docker.com/engine/install/ubuntu/}
        \item Docker Compose: \url{https://docs.docker.com/compose/install/linux/}
    \end{itemize}
    \item \inlinecode{sudo usermod -aG docker \$USER}を実行してください。
    \item Ubuntuの再起動をかけてください。
\end{enumerate}

\subsection{Visual Studio Code（VS Code）のインストール}
\begin{enumerate}
    \item Windows上にVS Code本体をインストールしてください。
    \item VS Code上の拡張機能を以下2つインストールしてください。偽の拡張機能（マルウェアの可能性があります）に注意してください。
    \begin{itemize}
        \item Dev Conteiners extension: \url{https://marketplace.visualstudio.com/items?itemName=ms-vscode-remote.remote-containers}
        \item Remote Development extension pack: \url{https://marketplace.visualstudio.com/items?itemName=ms-vscode-remote.vscode-remote-extensionpack}
    \end{itemize}
\end{enumerate}

\subsection{Gitの設定}
WindowsもしくはWSL2（Ubuntu）上にGitがインストールされているのは前提とします。ローカルにユーザー名とメールアドレスが適切に設定されていることを確認してください。
\begin{quote}
    \inlinecode{git config --global user.name "Your Name"}

    \inlinecode{git config --global user.email "your.email@address"}
\end{quote}

Dev Containerでは、適切に設定すれば、コンテナの内部でも上手くGitを使うことができます。つまり、ここで使う仮想コンテナは\inlinecode{./doc}に構築されるコンテナですが、そのコンテナはDigitShow Modbus本体のGitまで見に行ってくれます。よって、VS Codeでドキュメントを編集しながら、そのままpush等が可能です。ここではそのための設定を行います。

詳細はVS Codeのウェブサイト（\url{https://code.visualstudio.com/remote/advancedcontainers/sharing-git-credentials}）にある通りです。「Using a credential helper」というのが楽だと思います。これはつまりGitHubのウェブサイト（\url{https://docs.github.com/ja/get-started/getting-started-with-git/caching-your-github-credentials-in-git}）を参照せよということなので、これに従ってください。

\section{作業方法}
\begin{enumerate}
    \item VS Code上でDigitShow Modbusが格納されたディレクトリを開いてください。
    \item Ctrl + Shift + Pのショートカットを利用し、「開発コンテナー：コンテナーでフォルダを開く」を選択してください。
    \item \inlinecode{./doc}ディレクトリを選択しOKを押してください。
    \item 任意の\inlinecode{.tex}ファイルを編集し、保存すると自動でコンパイルが回ります。エディタ右上の▷を押してもコンパイルが回ります。▷の右の、分割画面にルーペがついたボタンを押すと、隣にPDFを参照しながら\inlinecode{.tex}ファイルを編集できます。
\end{enumerate}

\section{{\LaTeX}を書くときのルール}
\begin{itemize}
    \item 文書の見た目（スタイル）ではなく構造（意味や役割など）を意識してください。あくまで後者が主であり、前者はそれに付随する概念です。例えば、ある部分を太字にしたいと思っても、\inlinecode{{\textbackslash}textbf}は使わないでください。太字にしたいのはそこを強調するためであり、そうであれば強調コマンド\inlinecode{{\textbackslash}emph}を使うのが理にかなったやり方です。なお、\emph{強調はゴシック体}で定義されています。変更も可能です。
    \item ダブルクオーテーションマークで括るときは、\inlinecode{"}で括るのではなく\inlinecode{``}と\inlinecode{''}で括るようにするとダブルクオーテーションマークの形が整います。出力例は次の通りです。\inlinecode{"text."}{\rightarrow}"text." \inlinecode{``text.''}{\rightarrow}``text.''
    \item 好みの問題に近いですが、和文では全角記号を使うと見た目が整うように思います。例えば括弧類（）やコロン：など。和文の中に英数字を含めなければいけない場面では臨機応変に対応しましょう。
    \item 図表や章にはラベルをつけて参照することができます。ラベルのネーミングは（技術的には）自由ですが、\inlinecode{fig:}, \inlinecode{tab:}, \inlinecode{sec:}のように、何のラベルかを冒頭で示すと便利です。参照する際には\inlinecode{{\textbackslash}cref}を使うと、例えば図を参照した際には自動的に「図 1.1」のように、表を参照した際には自動的に「表 1.1」のように出力してくれます（cleverefパッケージ）。
    \item 単位つきの数値を扱うときにはsiunitxパッケージの機能を利用すると、数値と単位との間のスペーシングなどのスタイルが整います。例えば\qty{100}{\kPa}は\inlinecode{{\textbackslash}qty\{100\}\{{\textbackslash}kPa\}}とします。あるいは\qty{1e5}{\kg.\m^{-1}.\s^{-2}}は\inlinecode{{\textbackslash}qty\{1e5\}\{{\textbackslash}kg.{\textbackslash}m{\textasciicircum}\{-1\}.{\textbackslash}s{\textasciicircum}\{-2\}\}}とします。ほかにも\qtyrange{15}{20}{\kPa}は\inlinecode{{\textbackslash}qtyrange\{15\}\{20\}\{{\textbackslash}kPa\}}と書けます。
    \item 数式を書くときは、イタリック体と立体（ローマン体）の使い分けに注意しましょう。例えばイタリック体で$ABC$と書くと$A\times B\times C$と解釈されますが、立体で$\mathrm{ABC}$と書くとABC（一つの数）と解釈されます。また、添え字も$h_\mathrm{init}$にようにして立体にすることが多々あるようです。数式モードで立体にするには\inlinecode{{\textbackslash}mathrm\{\}}を使います。
    \item TikZパッケージを使うときれいな図が描けます。使い方はオンラインマニュアル（\url{https://tikz.dev/}）を見てください。
    \item コードやファイル名、ディレクトリをインラインで書きたいときは\inlinecode{{\textbackslash}inlinecode}を定義したので、そちらを使ってください。
    \item コードブロックもかけます（listingsパッケージ）。
    \begin{tcbcodeblock}{DigitShowModbus.cpp}
        \begin{lstcodeblock}[language=C++]
// DigitShowModbus.cpp : アプリケーション用クラスの機能定義を行います。
//

#include "pch.h"
#include "framework.h"
#include "afxwinappex.h"
#include "afxdialogex.h"
#include "DigitShowModbus.h"
#include "MainFrm.h"

#include "DigitShowModbusDoc.h"
#include "DigitShowModbusView.h"

#define _USE_MATH_DEFINES
#include <math.h>
#include <fstream>
#include <filesystem>
#include <cstdlib>

#include "sqlite3.h"
#include <spdlog/sinks/rotating_file_sink.h>
#include <nlohmann/json.hpp>
#include <modbus.h>

#ifdef _DEBUG
#define new DEBUG_NEW
#endif
        \end{lstcodeblock}
    \end{tcbcodeblock}
\end{itemize}
